\section{What a rate bound does and does not imply for learning}
\label{sec:critical-learning-comparison}

The controlled-circuit theorem is a target representation bound. A
learning theorem additionally needs to compare the inaccurate extensions
of a training sample with its true target. The next proposition shows
that normalized excess rate alone cannot supply a general complexity
measure for learning.

\begin{proposition}[Normalized action alone does not control learning]
\label{prop:critical-action-learning-obstruction}
There are depths $L_d=\Theta((d+1)^2)$ and a universal constant $C$
such that every Boolean function satisfies
\[
 \max\{0,L_d[J_{L_d}(f)-V_*]\}\le C(d+1)^2.
\]
Nevertheless, let $m=2^d$ and $n\le m/4$. For every learning rule
that assigns probabilities to the two labels after seeing $n$ uniform
cube inputs and their correct labels, there is a Boolean target $f$
for which
\[
 \Pr_S\!\left\{
   \Pr_X\{p_S(f(X)\mid X)<2/3\}\ge1/4
 \right\}\ge1-e^{-m/24}.
\]
This holds for iid inputs and for sampling without replacement.
Randomness of the learning rule may be included in the probability.
Consequently no universal sample bound can be polynomial solely in
$d$, $L$, and positive normalized excess $L[J_L(f)-V_*]$.
The hard target in this statement need not have a small circuit.
\end{proposition}

\begin{proof}
The uniform all-function rate upper bound in
Theorem~\ref{thm:critical-rate-flattening} gives
\[
 J_L(f)\le V_*+C\frac{(d+1)^2+\log^2(2L)}L
 \qquad\text{when }L\ge C_0(d+1).
\]
Take $L_d$ to be a sufficiently large universal multiple of
$(d+1)^2$, rounded up.

For the learning claim, first choose the entire target truth table
using independent fair signs. Condition on the sample inputs, their
observed labels, and the learner's randomness. At least
$u\ge m-n\ge3m/4$ inputs are unobserved. Their labels are still
independent fair signs, while the prediction probabilities are fixed.
At any such input at most one label can have probability at least
$2/3$, so the failure indicator has probability at least $1/2$.
These indicators are independent across the unobserved inputs.
Hoeffding's inequality gives
\[
 \Pr\{\text{fewer than }u/3\text{ unseen inputs fail}\}
 \le e^{-u/18}\le e^{-m/24}.
\]
Otherwise at least $u/3\ge m/4$ inputs fail. Averaging over the
random truth table supplies a deterministic target with the asserted
sample probability. Repeated observations only increase the number
of unobserved inputs. Finally, any fixed polynomial in the displayed
complexity parameters is eventually smaller than $2^d/4$.
The obstruction also applies to guarantees valid at all sufficiently
large widths with a target-dependent width threshold: at each fixed
$d$ there are finitely many targets, so one can take the maximum of
their thresholds before applying the learning lower bound.
\end{proof}

\paragraph{Why direct prior-mass comparisons retain a width cost.}
Write $\pi_N$ for the original globally margin-conditioned critical
prior on Boolean functions, and $J_{\min}=\min_gJ_L(g)$. At fixed
domain and depth the exact finite-class rates imply
\[
 -\frac1N\log\pi_N(f)\longrightarrow J_L(f)-J_{\min}.
\]
Thus a point-mass PAC-Bayes comparator retains a linear width cost
whenever the target rate exceeds the minimum. More generally, if a
comparator law $Q_N$ is supported on a nonempty function class
$\mathcal B$, then
\[
 D(Q_N\Vert\pi_N)\ge-\log\pi_N(\mathcal B)
 =N\left[\min_{g\in\mathcal B}J_L(g)-J_{\min}\right]+o(N).
\]
The inequality follows by conditioning $\pi_N$ on $\mathcal B$.
This does not refute width-independent learning; it explains why the
fixed-depth proofs instead control competing functions whose rates
are no larger than the target's.

\begin{proposition}[A target-relative error-witness criterion]
\label{prop:critical-error-witnesses}
Fix a Boolean target $f$, a distribution $\mathcal D$, and
$\epsilon,\delta\in(0,1)$. Define
\[
 \mathcal A_f=\{g:J_L(g)\le J_L(f)\}.
\]
Suppose there are sets $B_1,\ldots,B_M$ such that every
$g\in\mathcal A_f$ of error greater than $\epsilon/6$ has some
index $j$ with
\[
 B_j\subseteq\{x:g(x)\ne f(x)\},\qquad
 \mathcal D(B_j)\ge\epsilon/12.
\]
The sets may depend on $f$ and $\mathcal D$, but not on the sample.
If
\[
 n\ge\frac{12}{\epsilon}
       \left[\log M+\log(1/\delta)\right],
\]
then, at all sufficiently large learner widths, the original critical
posterior has confidence error
$\mathcal D\{x:p_N(f(x)\mid x)<2/3\}\le\epsilon$
with sample probability at least $1-\delta$.
The width threshold can be chosen uniformly over the realized sample.
\end{proposition}

\begin{proof}
An inaccurate consistent function in $\mathcal A_f$ forces the sample
to miss one of the $B_j$. A union bound bounds this event by
$M\exp(-n\epsilon/12)\le\delta$. On the complementary event,
all consistent members of $\mathcal A_f$ have error at most
$\epsilon/6$.

Finiteness of the Boolean function class and the exact prior rates give
\[
 \frac{\pi_N(\mathcal A_f^c)}{\pi_N(f)}\longrightarrow0.
\]
For every sample labeled by $f$, its posterior normalizer is at least
$\pi_N(f)$. Thus, for all sufficiently large widths, posterior mass
outside $\mathcal A_f$ is at most $\epsilon/6$, uniformly over
the sample. On the good-sample event its Gibbs error is at most
$\epsilon/3$. At every point of confidence error the posterior
wrong-label probability exceeds $1/3$, giving the claimed bound.
Mixtures of tied rate minimizers are included in this argument.
\end{proof}

The witness condition is a one-sided containment statement about error
sets. An ordinary $L^1(\mathcal D)$ cover does not suffice: a predictor
can change its labels on a small-measure set chosen to contain the
observed points. Taking one witness for each inaccurate function
recovers the usual cardinality proof. A substantially smaller common
family could give a learning theorem even when that proof is wasteful;
the controlled-circuit rate upper bound alone does not provide such
a family.
